% Work-MAT | 2. rocnik | pripravovana hodina 27
% Návaznost: uživatelské připravy.zip (hodiny 01–20), autorská učebnice na GitHubu.
% Realisticky.cz: didaktická inspirace v posloupnosti a typech úloh; zadání jsou vlastní.
\documentclass[11pt,a4paper]{article}
\input{../../sablony/pripravy/priprava.sty}
\nastavPripravu{2. ročník}{Soustavy exponenciálních rovnic}{27}{Soustavy exponenciálních rovnic II}
\begin{document}
\blok{Motivace a návaznost}
V předchozí hodině měly pomocné proměnné rovnou podobu lineární soustavy. Dnes nejprve upravíme exponenty a koeficienty.
\poznamka{Orientační nabídka pro 45--55 minut: krátké zopakování a motivace, nový výklad, tři společné vzorové úlohy, samostatná práce. Vyber podle tempa třídy; zbývající zadání slouží jako rezerva.}
\blok{Výklad}
Dvojice $u=2^x$ a $v=3^y$ umožní řešit smíšené soustavy po běžné algebraické úpravě. Přijímáme jen kladná pomocná řešení.
\blok{Společné příklady}
\begin{prikladReseniBox}[title={Příklad 1}]
Vyřeš, resp. urči: $\begin{cases}2^x+3^y=11\\2^{x+1}-3^y=13\end{cases}$.
\tcblower
\[\begin{aligned}u=2^x>0,\quad v=3^y>0 \\ \begin{cases}u+v=11\\2u-v=13\end{cases} \\ u=8,\;v=3\;\Longrightarrow\;(x,y)=(3,1)\end{aligned}\]
\end{prikladReseniBox}
\begin{prikladReseniBox}[title={Příklad 2}]
Vyřeš, resp. urči: $\begin{cases}2^x+2^y=10\\2^{x+1}-2^y=14\end{cases}$.
\tcblower
\[\begin{aligned}u+v=10,\quad 2u-v=14 \\ u=8,\quad v=2 \\ (x,y)=(3,1)\end{aligned}\]
\end{prikladReseniBox}
\poznamka{Při rozboru ukaž, že soustava může mít dvě rovnice v různých základech, ale stále jen dvě pomocné neznámé.}
\blok{Samostatné procvičování}
\begin{zadaniReseniBox}[title={Samostatná práce 1}]
\begin{enumerate}[label=\alph*)]
\item $2^x+3^y=11,\;2^{x+1}-3^y=13$.
\item $2^x+2^y=10,\;2^{x+1}-2^y=14$.
\item $3^x+2^y=13,\;3^x-2^y=5$.
\item Jak poznáš nepřípustné pomocné řešení?.
\end{enumerate}
\tcblower
\begin{enumerate}[label=\alph*)]
\item $(3,1)$.
\item $(3,1)$.
\item $(2,2)$.
\item $u\le0$ nebo $v\le0$..
\end{enumerate}
\end{zadaniReseniBox}
\blok{Dobrovolné rozšíření}
\blok{Úkol na doma}
Při početních úlohách napiš postup, zvolenou metodu a zkoušku v původním zadání. U soustav kontroluj obě rovnice; u nerovnic výsledek zakresli na osu.
\begin{zadaniReseniBox}[title={Domácí úkol}]
\begin{enumerate}[label=\alph*)]
\item $2^x+3^y=11,\;2^{x+1}-3^y=13$.
\item $2^x+2^y=10,\;2^{x+1}-2^y=14$.
\item $3^x+2^y=13,\;3^x-2^y=5$.
\item $2^x+2^y=6,\;2^{x+1}-2^y=6$.
\item $2^x+3^y=7,\;2^x-3^y=1$.
\item Proč ověřit obě rovnice?.
\end{enumerate}
\tcblower
\begin{enumerate}[label=\alph*)]
\item $(3,1)$.
\item $(3,1)$.
\item $(2,2)$.
\item $(2,1)$.
\item $(2,1)$.
\item $Kořen musí vyhovět celé soustavě.$.
\end{enumerate}
\end{zadaniReseniBox}
\poznamka{Domácí úkol obsahuje 6 položek; předpokládaná domácí práce přibližně 15--30 minut. Po větších blocích jsou vědomě promíchané dříve probrané typy.}
\end{document}
